\documentclass[aspectratio=43]{beamer}
\usetheme{CCNU}

\usepackage{slashed}
\usepackage{datetime}
\usepackage{graphicx}
\usepackage{hyperref}
\usepackage{cancel}
\usepackage[absolute,overlay]{textpos}
\usepackage{amsmath,amssymb,booktabs}
\usepackage{listings}
\usepackage{xcolor}

% If your images are stored in a folder, update this path.
\graphicspath{{./}{plots/}}

\def\Put(#1,#2)#3{\leavevmode\makebox(0,0){\put(#1,#2){#3}}}
\def\mydate{\leavevmode\hbox{\the\year-\twodigits\month-\twodigits\day}}
\def\twodigits#1{\ifnum#1<10 0\fi\the#1}

\newcommand\Wider[2][3em]{%
\makebox[\linewidth][c]{%
  \begin{minipage}{\dimexpr\textwidth+#1\relax}
  \raggedright
  \centering#2
  \end{minipage}%
  }%
}

\def\meetingname{Predictive Analytics}

\title[\meetingname]{The Riemann Zeta Function: Zeros, the Critical Line, and Computational Exploration}
\author[G Shan]{Samuel Castillo}
\institute[UMass Amherst]{University of Massachusetts Amherst}
\date[\mydate]{\meetingname\\\today}

\lstset{
    basicstyle=\ttfamily\scriptsize,
    breaklines=true,
    frame=single,
    columns=fullflexible,
    showstringspaces=false,
    keywordstyle=\color{blue},
    commentstyle=\color{gray},
    stringstyle=\color{red}
}

\begin{document}

\begin{frame}[plain,t]
\titlepage
\end{frame}

\begin{frame}
\frametitle{Contents}
\tableofcontents
\end{frame}

%------------------------------------------------------------
\section{Introduction}

\begin{frame}{Introduction}
\begin{itemize}
    \item For $\Re(s)>1$, the Riemann zeta function is
    \[
    \zeta(s)=\sum_{n=1}^{\infty}\frac{1}{n^s}.
    \]
    \item Euler's product formula links the zeta function directly to the primes:
    \[
    \zeta(s)=\prod_p\left(1-p^{-s}\right)^{-1}.
    \]
    \item Central thesis: the non-trivial zeros of $\zeta(s)$ govern the oscillatory behavior of prime counting.
\end{itemize}
\end{frame}

\section{From Basel to Prime Number Theory}

\begin{frame}{From Euler's Basel Problem to Primes}
\begin{itemize}
    \item Euler solved the Basel problem:
    \[
    \sum_{n=1}^{\infty}\frac{1}{n^2}=\frac{\pi^2}{6},\qquad \zeta(2)=\frac{\pi^2}{6}.
    \]
    \item This showed that a purely arithmetic series could evaluate to an expression involving $\pi$.
    \item Riemann later extended $\zeta(s)$ to the complex plane (except for a simple pole at $s=1$) and connected its zeros to the distribution of primes.
\end{itemize}
\end{frame}

\begin{frame}{Prime Density and the Big Picture}
\begin{itemize}
    \item Prime Number Theorem:
    \[
    \pi(x)\sim \frac{x}{\log x}.
    \]
    \item The zeros of $\zeta(s)$ control deviations from this leading approximation.
    \item Related prime phenomena:
    \begin{itemize}
        \item Mersenne primes: $2^p-1$ with $p$ prime necessary but not sufficient.
        \item Twin primes: pairs $(p,p+2)$ such as $(29,31)$.
    \end{itemize}
\end{itemize}
\end{frame}

\section{Visualizing Zeta Zeros}

\begin{frame}{Visualizing Zeta Function Zeros}
\begin{itemize}
    \item Non-trivial zeros lie in the critical strip
    \[
    0<\Re(s)<1.
    \]
    \item The Riemann Hypothesis asserts that all non-trivial zeros lie on
    \[
    \Re(s)=\tfrac12.
    \]
    \item Three useful visualization modes:
    \begin{itemize}
        \item \textbf{Phase plots}: hue encodes $\arg(\zeta(s))$ and brightness encodes $|\zeta(s)|$.
        \item \textbf{Contour plots}: intersections of $\Re(\zeta(s))=0$ and $\Im(\zeta(s))=0$ locate zeros.
        \item \textbf{3D surfaces}: minima of $|\zeta(s)|$ reveal the zero set geometrically.
    \end{itemize}
\end{itemize}
\end{frame}

\begin{frame}{Phase Plot (Domain Coloring)}
\centering
% Update the filename below if your local file uses a different name.
\includegraphics[width=0.82\textwidth,height=0.66\textheight,keepaspectratio]{zeta_phase_corrected.png}

\vspace{0.4em}
{\small Hue represents $\arg(\zeta(s))$, brightness represents $|\zeta(s)|$, and zeros appear as phase singularities near $\Re(s)=\tfrac12$.}
\end{frame}

\begin{frame}{Three-Dimensional Surface of $|\zeta(s)|$}
\centering
% Update the filename below if your local file uses a different name.
\includegraphics[width=0.82\textwidth,height=0.66\textheight,keepaspectratio]{z3d.png}

\vspace{0.4em}
{\small The non-trivial zeros appear as localized minima of $|\zeta(s)|$ rather than as a continuous trench along the critical line.}
\end{frame}

\begin{frame}{Numerical Contour Plot in the Critical Strip}
\centering
% Update the filename below if your local file uses a different name.
\includegraphics[width=0.82\textwidth,height=0.66\textheight,keepaspectratio]{rzt.png}

\vspace{0.4em}
{\small Red curves show $\Re(\zeta(s))=0$, green curves show $\Im(\zeta(s))=0$, and their intersections align numerically with the critical line.}
\end{frame}

\section{The Critical Line and RH}

\begin{frame}{The Critical Line and the Riemann Hypothesis}
\begin{itemize}
    \item Functional equation:
    \[
    \zeta(s)=2^s\pi^{s-1}\sin\!\left(\frac{\pi s}{2}\right)\Gamma(1-s)\zeta(1-s).
    \]
    \item This relates $s$ and $1-s$ and extends $\zeta(s)$ to the complex plane apart from the pole at $s=1$.
    \item \textbf{Riemann Hypothesis:} every non-trivial zero lies on $\Re(s)=\tfrac12$.
\end{itemize}
\end{frame}

\begin{frame}{Why the Hypothesis Matters}
\begin{itemize}
    \item Extensive computation has verified enormous numbers of zeros on the critical line.
    \item The zeros enter prime-counting formulas such as
    \[
    \pi(x)=\mathrm{Li}(x)-\sum_\rho \mathrm{Li}(x^\rho)+\text{lower-order terms}.
    \]
    \item RH implies an optimal-order error term in the Prime Number Theorem:
    \[
    \pi(x)=\mathrm{Li}(x)+O\!\left(x^{1/2}\log x\right).
    \]
\end{itemize}
\end{frame}

\section{Explicit Formula and Prime Density}

\begin{frame}{Prime Density and the Explicit Formula}
\begin{itemize}
    \item A central bridge between zeros and primes is the explicit formula
    \[
    \psi(x)=x-\sum_\rho \frac{x^\rho}{\rho}-\log(2\pi)-\frac12\log(1-x^{-2}).
    \]
    \item The zeros contribute oscillatory corrections to the main term $x$.
    \item Related asymptotics:
    \[
    P(n\text{ prime})\approx \frac{1}{\ln n}, \qquad \mathrm{Li}(x)\sim \frac{x}{\ln x}.
    \]
\end{itemize}
\end{frame}

\section{Computational and AI Methods}

\begin{frame}{Computational and AI-Assisted Approaches}
\begin{itemize}
    \item \textbf{Numerical contour analysis:} approximate zeros and visualize local spacing behavior.
    \item \textbf{Automated theorem search:} explore large proof spaces and useful intermediate lemmas.
    \item \textbf{Large-scale verification:} test zero locations to increasing heights.
    \item These methods do \emph{not} prove RH, but they expand the range of meaningful experiments.
\end{itemize}
\end{frame}

\section{Cryptographic Motivation}

\begin{frame}{Cryptographic Motivation}
\begin{itemize}
    \item Classical ciphers (Caesar shifts, keyword substitution) are broken by brute force or frequency analysis.
    \item Modern public-key systems rely on hard number-theoretic problems such as integer factorization.
\end{itemize}

\vspace{0.6em}
\centering
\begin{tabular}{lccc}
\toprule
System & Key space / scale & Attack & Break time \\
\midrule
Caesar cipher & 25 keys & Brute force & Instant \\
Keyword substitution & structured & Frequency analysis & Seconds--minutes \\
DES (56-bit) & $2^{56}$ & Exhaustive search & Hours--days \\
RSA-2048 & very large modulus & Factorization & Infeasible classically \\
\bottomrule
\end{tabular}
\end{frame}

\section{Conclusion and Future Work}

\begin{frame}{Conclusion}
\begin{itemize}
    \item Euler opened one door into the zeta function through special values such as $\zeta(2)=\pi^2/6$.
    \item Riemann opened a deeper door: the zeros of $\zeta(s)$ govern the fine structure of the primes.
    \item Computational plots make this connection concrete, even though they do not constitute a proof.
\end{itemize}
\end{frame}

\begin{frame}{Future Work}
\begin{itemize}
    \item High-precision numerical approximation of $\zeta(s)$ along the critical line.
    \item Improved visualization using domain coloring, contour plots, and surface plots.
    \item Computational experiments with truncated explicit formulas for $\psi(x)$ and $\pi(x)$.
    \item Exploratory AI-based analysis of zero spacing, clustering, and prime-related data.
\end{itemize}
\end{frame}

\appendix
\section{Appendix: Reproducibility}

\begin{frame}[fragile]{Appendix: Python Code for the Phase Plot}
\begin{lstlisting}[language=Python]
import numpy as np
import matplotlib.pyplot as plt
import mpmath as mp
from matplotlib.colors import hsv_to_rgb

mp.dps = 30
re_vals = np.linspace(0.35, 0.65, 500)
im_vals = np.linspace(10.0, 32.0, 700)
Re, Im = np.meshgrid(re_vals, im_vals)
phase = np.zeros_like(Re, dtype=float)
magnitude = np.zeros_like(Re, dtype=float)

for i in range(Re.shape[0]):
    for j in range(Re.shape[1]):
        s = mp.mpc(float(Re[i, j]), float(Im[i, j]))
        z = mp.zeta(s)
        phase[i, j] = float(mp.arg(z))
        magnitude[i, j] = float(abs(z))

phase_norm = (phase + np.pi) / (2 * np.pi)
value = 1.0 / (1.0 + magnitude)
hsv = np.dstack((phase_norm, np.ones_like(value), value))
rgb = hsv_to_rgb(hsv)
\end{lstlisting}
\end{frame}

\begin{frame}[fragile]{Appendix: Python Code for the Surface Plot}
\begin{lstlisting}[language=Python]
import numpy as np
import matplotlib.pyplot as plt
import mpmath as mp
from mpl_toolkits.mplot3d import Axes3D

mp.dps = 30
re_vals = np.linspace(0.2, 0.8, 160)
im_vals = np.linspace(10.0, 35.0, 220)
Re, Im = np.meshgrid(re_vals, im_vals)
Z_abs = np.zeros_like(Re, dtype=float)

for i in range(Re.shape[0]):
    for j in range(Re.shape[1]):
        s = mp.mpc(float(Re[i, j]), float(Im[i, j]))
        Z_abs[i, j] = float(abs(mp.zeta(s)))

Z_plot = np.log1p(Z_abs)
# plot_surface(Re, Im, Z_plot) and save figure as z3d.png
\end{lstlisting}
\end{frame}

\setbeamercolor{background canvas}{bg=CCNUMaize}
\begin{frame}[plain,t]
\vspace{100pt}
\centering
\Large Thank you
\end{frame}

\end{document}
